\section{Problem Setting and Background}
\label{sec:background}

\subsection{Reward-free pretraining and successor measures}

% Define the reward-free MDP and successor measure used throughout the paper.
We consider a reward-free Markov decision process $(\Sspace,\Aspace,\Ptrans,\gamma)$ with state space $\Sspace$, action space $\Aspace$, transition kernel $\Ptrans$, and discount factor $\gamma$. For a policy $\pi$, the successor measure from $(s,a)$ to a measurable set $X\subseteq \Sspace$ is
\begin{equation}
    M^{\pi}(s,a,X)
    =
    \E_{\pi}\left[
        \sum_{t\geq 0}\gamma^t\ind\{s_{t+1}\in X\}
        \mid s_0=s,a_0=a
    \right].
\end{equation}
When the state space is discrete, we write $M^{\pi}(s,a,s^+)$ for the discounted occupancy of state $s^+$.

% Successor measures turn reward-conditioned policy evaluation into a linear objective.
For a state reward $r$, policy evaluation can be written as
\begin{equation}
    Q^{\pi}_r(s,a)
    =
    \int_{\Sspace} M^{\pi}(s,a,ds^+) r(s^+).
\end{equation}
This identity explains why successor measures are useful for URL: a reward-free object that predicts future occupancy can be combined with a downstream reward to evaluate or infer policies \citep{touati-neurips21a,touati-iclr23a,agarwal-rlj25a}.

% Practical methods compress the successor object; this creates the relevant risk.
The exact successor measure is intractable to represent for a large policy class. Practical URL methods therefore impose parametric structure. \fb{} uses a low-rank approximation,
\begin{equation}
    M^{\pi_z}(s,a,s^+)
    \approx
    F(s,a,z)^{\top}B(s^+),
\end{equation}
where $z$ indexes a reward or policy embedding \citep{touati-neurips21a,touati-iclr23a}. Goal-conditioned and contrastive methods use target-conditioned values, policies, or embeddings instead of explicitly representing the full measure \citep{eysenbach-iclr21a,eysenbach-neurips22a,ma-neurips22a}. \psm{} and related basis methods approximate successor or value objects through learned bases \citep{mahadevan-icml05a,mahadevan-jmlr07a,farebrother-iclr23a,agarwal-arxiv25a}. In each case, the representation must decide which states are equivalent for downstream control.

\subsection{Manipulation as phase-conditioned reachability}

% ontact phase changes the transition kernel.
We model a manipulation state as \(s=(x,o,g,c),\) where $x$ denotes robot proprioception, $o$ denotes object pose, $g$ denotes goal pose, and $c$ denotes a contact phase.
The phase variable may be latent in the observation, but it changes how actions affect the object.
We write the transition kernel as
\[
      P(s'\mid s,a) = P_c(x',o',c'\mid x,o,c,a).
\]
For cube manipulation, the relevant phases include free approach, contact acquisition, stable grasped transport, unstable contact, and dropped contact.
% Geometric proximity and controllability are different.
This formulation separates geometric proximity from controllability.
When the cube is not grasped, actions that move the gripper need not move the cube.
When the cube is stably grasped, similar gripper actions transport the cube.
When contact is lost, the same object-goal distance can correspond to a different set of feasible future trajectories.
As a result, a representation that preserves only object-goal geometry may merge states whose future occupancy of successful transport states is different.
